%!TEX root = ../template.tex
%%%%%%%%%%%%%%%%%%%%%%%%%%%%%%%%%%%%%%%%%%%%%%%%%%%%%%%%%%%%%%%%%%%%
%% chapter-manual-defaults.tex
%% NOVA thesis document file
%%
%% Manual for creating and customizing School Default Files (.ldf)
%%%%%%%%%%%%%%%%%%%%%%%%%%%%%%%%%%%%%%%%%%%%%%%%%%%%%%%%%%%%%%%%%%%%

\typeout{NT FILE chapter-manual-defaults.tex}%

\chapter{Creating School Configuration Files}
\label{cha:defaults_files}

The \gls{novathesis} template uses configuration files (with the \texttt{.ldf} extension) to define the specific rules, strings, and layouts for each supported University and School. These files are located in the \texttt{NOVAthesisFiles/Schools/} directory and are loaded based on the \texttt{school} option in \texttt{0-Config/1\_novathesis.tex}.

\section{Identification and Metadata}
\label{sec:identification}

The first part of a defaults file typically defines the names and logos for the University and School.

\subsection{University and School Names}
Use the \latexinline{\university} and \latexinline{\school} commands to define localized names. These commands use the \texttt{MemStore} system, allowing you to associate values with language codes.

\begin{ltx}
% University
\university(pt):={Universidade NOVA de Lisboa}
\university(en):={NOVA University Lisbon}

% Faculty / School
\school(pt):={Faculdade de Ciências e Tecnologia}
\school(en):={NOVA School of Science and Technology}
\end{ltx}

\subsection{Logos}
Logos are also stored in memory and can have multiple variants (e.g., color vs. monochrome, positive vs. negative).

\begin{ltx}
\school(logo,pt):={nova-fct-logo-pt}
\school(logo,neg,pt):={nova-fct-logo-negative-pt}

% Dynamic selection based on the cover language
\school(logo):={\theschool(logo,\@LANG@COVER)}
\end{ltx}

\section{Colors}
\label{sec:colors}

Defaults files often define the official colors of the University or School using the \texttt{xcolor} package. These colors can then be used in cover elements, hooks, or spine configurations.

\begin{ltx}
\definecolor{iselRed}{HTML}{FF3C28}
\definecolor{iselPlum}{HTML}{5F1437}
\end{ltx}

\section{Document Configuration}
\label{sec:config}

\subsection{Template Options}
You can use \latexinline{\ntsetup} inside a defaults file to enforce specific settings for a school.

\begin{ltx}
\ntsetup{print/frontpage=true}
\ntsetup{style/font=verdana}
\end{ltx}

\subsection{Committee Order}
The order in which committee members are printed is defined by \latexinline{\committeeorder}. The keys are: \texttt{c} (chair), \texttt{r} (rapporteur), \texttt{a} (adviser), \texttt{ca} (co-adviser), \texttt{m} (member), \texttt{g} (guest).

\begin{ltx}
\committeeorder():={c,r,a,m}%
\end{ltx}

\section{Customizing Strings and Templates}
\label{sec:strings}

\subsection{Dissertation Strings}
The text that appears on the cover describing the purpose of the document (e.g., "Dissertation to obtain the degree of...") is defined via \latexinline{\dissertationstring}. It can be specialized by document type and language.

\begin{ltx}
\dissertationstring(msc,en) = {%
  \GetDocType(\option{/novathesis/doctype}) to obtain the degree of %
  \thedegreestring(\option{/novathesis/doctype},\expanded{\thedocauthor(gender)},\@LANG@COVER)%
  \ in \themajorfield(\@LANG@COVER).%
}
\end{ltx}

\subsection{String Templates}
The template provides several commands to customize the strings used for advisers, committee members, and degrees. These are usually stored in memory and can be redefined per language and gender.

\begin{ltx}
% Redefining adviser strings for English
\adviserstring(a,1,m,en):={Supervisor}
\adviserstring(a,1,f,en):={Supervisor}
\adviserstring(c,1,m,en):={Co-supervisor}
\end{ltx}

\subsection{New Data Fields}
You can create custom fields using \latexinline{\newdata} and set them using the command name.

\begin{ltx}
\newdata{provisionalstring}
\provisionalstring(pt):={Documento Provisório}
\provisionalstring(en):={Provisional Document}
\end{ltx}

\section{Copyright and Legal Strings}
\label{sec:copyright}

The legal text for the copyright page can be customized using \latexinline{\copyrighttextstring}.

\begin{ltx}
\copyrighttextstring(pt)={A \theschool(\@LANG@COPYRIGHT) e a \theuniversity(\@LANG@COPYRIGHT) têm o direito perpétuo de arquivar este relatório...}
\end{ltx}

\section{School-specific Statements}
\label{sec:statements}

Many schools require specific "Statements of Integrity" or "Author's Statements." You can register these files using \latexinline{\ntaddfile}.

\begin{ltx}
\ntaddfile{statement}[en]{ipl-isel/ipl-isel-statement-en}
\ntaddfile{statement}[pt]{ipl-isel/ipl-isel-statement-pt}
\ntsetup{print/statement=true}
\end{ltx}

\section{Batch Definitions with Loops}
\label{sec:loops}

To avoid repetitive code when defining strings for multiple languages or variants, you can use the \latexinline{\@for} loop (or the template's internal \latexinline{\ntfor}).

\begin{ltx}
\@for\mylang:={\@LANG@LOADED}\do{%
  \committeetitlestring(\mylang):={\ToUppercase{\thecommitteetitlestring(\mylang)}}%
}
\end{ltx}

\section{Layout and Margins}
\label{sec:margins}

Margins are defined for different media (\texttt{screen}, \texttt{paper}) and specific pages (\texttt{cover}).

\begin{ltx}
\margin(screen,top):={2.5cm}
\margin(paper,left):={3.5cm}
\margin(cover,top):={2.5cm}
\end{ltx}

\section{Spine Configuration}
\label{sec:spine}

The book spine layout is managed by the \latexinline{\spine} command and the \latexinline{\SpineSetup} helper.

\begin{ltx}
\spine(logo)={ipl-isel-logo/90}
\spine(author)={\thedocauthor(name,short)}

\SpineSetup{
  bg/image = {my-spine-background},
  box/title/fg/color = {white},
  box/logo/len = {2.1cm},
}
\end{ltx}

\section{Customizing Covers with \textbackslash ntaddtocover}
\label{sec:ntaddtocover}

The \latexinline{\ntaddtocover} command is the primary tool for placing elements on the cover pages (\texttt{1-1}, \texttt{2-1}, etc.).

\subsection{Command Syntax}
\begin{center}
  \latexinline{\ntaddtocover[<attributes>]{<covers>}{<code>}}
\end{center}

\subsection{Available Attributes}
\begin{description}
  \item[\texttt{vspace}] Vertical spacing. Supports lengths (\texttt{1cm}) or integers (relative \latexinline{\vfill} weight).
  \item[\texttt{hspace}] Horizontal spacing.
  \item[\texttt{height} / \texttt{width}] Box dimensions.
  \item[\texttt{halign} / \texttt{valign}] Alignment (\texttt{l, c, r, j} and \texttt{t, c, b}).
  \item[\texttt{status} / \texttt{degree}] Conditional filters.
  \item[\texttt{tikz}] If \texttt{true}, code is executed in a page-anchored TikZ environment.
  \item[\texttt{color} / \texttt{bgcolor}] Text and background colors for the container box.
\end{description}

\subsection{Resetting Covers}
Use \latexinline{\ntresetcovers} to clear all previously defined cover elements.

\section{Debugging and Reference Layouts}
\label{sec:debugging}

When designing a new cover layout, it is often necessary to align elements precisely according to school regulations or a provided reference PDF.

\subsection{Millimeter Grid (Debug Mode)}
You can enable a millimeter grid that covers the entire printing area of the covers and the spine by adding \texttt{debug=cover} to your \latexinline{\ntsetup} in \texttt{0-Config/1\_novathesis.tex}.

\begin{ltx}
\ntsetup{debug=cover}
\end{ltx}

This will render a semi-transparent grid with 1mm, 5mm, 1cm, and 5cm steps, helping you measure distances and verify alignments.

\subsection{Using a Reference Background Image}
A common technique for replicating an existing cover design is to load a scan or a PDF of the official cover into the background, and then place your \latexinline{\ntaddtocover} elements on top of it.

\begin{enumerate}
  \item \textbf{Load the reference image}: In your defaults file, set the background image for the desired cover page (usually \texttt{1-1} or \texttt{2-1}).
  \begin{ltx}
  % Load reference-cover.pdf as the background
  \thesiscover(1-1,image):={reference-cover}
  \end{ltx}
  \item \textbf{Align your elements}: Use \latexinline{\ntaddtocover} with attributes like \texttt{vspace}, \texttt{hspace}, \texttt{height}, and \texttt{halign} to position your text and logos.
  \begin{ltx}
  \ntaddtocover[vspace=45mm, halign=l]{1-1}{%
    \fontsize{20}{24}\selectfont\textbf{My Dissertation Title}
  }
  \end{ltx}
  \item \textbf{Finalize}: Once all elements are correctly positioned, remove the reference image line from your configuration and turn off the debug flag.
\end{enumerate}

\section{Hooks}
\label{sec:hooks}

Hooks allow you to inject code at specific points in the rendering process.

\begin{ltx}
\NTAddToHook{cover/1-1/text/pre}{%
  \sffamily\color{mySchoolColor}%
}
\end{ltx}

\section{Tips \& Tricks}
\label{sec:tips}

\begin{itemize}
  \item \textbf{Expansion}: Use \latexinline{\expanded{...}} when passing memory data (like \latexinline{\theschool(...)}) to commands that don't automatically expand their arguments, such as \latexinline{\includegraphics}.
  \item \textbf{Conditionals}: Use \latexinline{\ifphddoc}, \latexinline{\ifmscdoc}, \latexinline{\ifdocfinal}, etc., to tailor the configuration to the current document state.
  \item \textbf{TikZ Coordinate System}: When using \latexinline{[tikz]}, the \texttt{(current page)} node is available. The Y-axis is often inverted in these environments, so check school examples.
  \item \textbf{Inheritance}: Defaults files often \latexinline{\input} other defaults files. Use \latexinline{\ntresetcovers} if you want to completely replace the cover layout of the parent file.
  \item \textbf{Debug Mode}: Use \latexinline{\ntsetup{debug=cover}} to see the bounding boxes of your cover elements.
  \item \textbf{Case-Insensitive Match}: Use \latexinline{\datamatch} for a flexible lookup across multiple possible memory keys.
  \item \textbf{Utility Commands}: You can define small helper commands in defaults files to simplify layout logic (e.g., \latexinline{\printdwithvspace}).
  \item \textbf{Loading Fonts}: Use \latexinline{\AtEndPreamble} to load school-specific fonts or packages to avoid clashes with the class's early loading sequence.
\end{itemize}
